\documentclass[11pt]{article}

\usepackage[utf8]{inputenc}
\usepackage[T1]{fontenc}
\usepackage{amsmath,amssymb,amsthm}
\usepackage{algorithm}
\usepackage{algpseudocode}
\usepackage{booktabs}
\usepackage{graphicx}
\usepackage{hyperref}
\usepackage{xcolor}
\usepackage{listings}
\usepackage[margin=1in]{geometry}
\usepackage{enumitem}
\usepackage{url}
\usepackage{textcomp}
\usepackage{tikz}
\usepackage{multirow}
\usepackage{caption}
\usepackage{subcaption}
\usepackage{tabularx}

\newtheorem{theorem}{Theorem}
\newtheorem{lemma}[theorem]{Lemma}
\newtheorem{corollary}[theorem]{Corollary}
\newtheorem{proposition}[theorem]{Proposition}
\newtheorem{definition}[theorem]{Definition}
\newtheorem{remark}[theorem]{Remark}

\newcommand{\bioprover}{\textsc{BioProver}}
\newcommand{\cegar}{\textsc{CEGAR}}
\newcommand{\cegis}{\textsc{CEGIS}}
\newcommand{\dreach}{\textsc{dReach}}
\newcommand{\breach}{\textsc{Breach}}
\newcommand{\dreal}{\textsc{dReal}}
\newcommand{\stl}{\textsc{STL}}
\newcommand{\biostl}{\textsc{Bio-STL}}
\newcommand{\RR}{\mathbb{R}}
\newcommand{\NN}{\mathbb{N}}

\lstset{
  language=Python,
  basicstyle=\ttfamily\small,
  keywordstyle=\color{blue!70!black},
  commentstyle=\color{green!50!black},
  stringstyle=\color{red!60!black},
  breaklines=true,
  frame=single,
  columns=fullflexible,
  captionpos=b,
}

\title{\bioprover{}: AI-Guided CEGAR for Formal Verification and\\
Repair of Synthetic Biology Circuits}

\date{}

\begin{document}
\maketitle

\begin{abstract}
Synthetic biology circuits are designed by composing genetic parts
(promoters, repressors, activators) into regulatory networks that
implement logic, oscillation, and memory functions. Verifying that a
circuit design satisfies its temporal specification---given uncertain
kinetic parameters---remains an open challenge: the dynamics are
nonlinear, high-dimensional, and stiff. We present \bioprover{}, a
counterexample-guided abstraction refinement (\cegar{}) framework
that formally verifies gene circuit designs against
signal temporal logic (\stl{}) specifications over ODE models. Key
contributions include: (1)~a biology-specific refinement strategy
exploiting Hill-function monotonicity that reduces \cegar{}
convergence from exponential to $O(m \log(x_{\max}/\varepsilon))$
iterations for monotone circuits; (2)~an explicit four-level
soundness hierarchy (\textsc{Sound}, \textsc{$\delta$-Sound},
\textsc{Bounded}, \textsc{Approximate}) with end-to-end error
tracking through $\delta$-decidability, discretization, and moment
closure; (3)~compositional verification via circular
assume-guarantee reasoning with formalized well-formedness
conditions; and (4)~an optional ML-accelerated predicate predictor
that speeds up refinement by $1.54\times$ while never compromising
soundness. On a benchmark suite of 29 circuits (2--20 species),
\bioprover{} verifies 15 circuits with full sound guarantees, provides
bounded guarantees for 13, and scales compositionally to 15 species
where monolithic verification times out.
\end{abstract}

% ===================================================================
\section{Introduction}
\label{sec:intro}
% ===================================================================

Synthetic biology aims to engineer living cells with programmable
behavior by assembling genetic circuits from standardized
parts~\cite{endy2005foundations}. A toggle switch combining two
mutually repressing genes produces bistable memory; a repressilator
with three genes in a negative-feedback ring generates sustained
oscillations~\cite{elowitz2000synthetic}. Before fabricating these
circuits in the wet lab, designers need assurance that the intended
behavior holds across the range of plausible kinetic parameters
obtained from part characterization data.

The verification challenge is formidable. Gene circuit dynamics are
governed by systems of nonlinear ordinary differential equations
(ODEs) with Hill-function kinetics:
\begin{equation}
\label{eq:hill}
  \dot{x}_i = \sum_j V_{ij}\,\frac{x_j^{n_{ij}}}{K_{ij}^{n_{ij}} + x_j^{n_{ij}}}
  - \gamma_i\,x_i + u_i,
\end{equation}
where parameters $(V_{ij}, K_{ij}, n_{ij}, \gamma_i)$ are uncertain.
Specifications such as ``protein~$X$ always remains above
concentration~0.5'' or ``the circuit oscillates with period between
18 and 25 time units'' are naturally expressed in signal temporal
logic (\stl{})~\cite{maler2004monitoring}.

Existing approaches fall into three categories.
\emph{Simulation-based} tools like \breach{}~\cite{donze2010breach}
sample parameter space and check \stl{} robustness but provide no
formal guarantees.
\emph{Reachability tools} like Flow*~\cite{chen2013flow} and
\dreach{}~\cite{kong2015dreach} compute flowpipe
over-approximations but do not natively support \stl{} specifications
or parameter synthesis.
\emph{Statistical model checking}~\cite{younes2006statistical}
provides probabilistic guarantees but cannot prove universal
properties. None of these approaches combine formal verification with
automatic parameter repair.

\paragraph{Contributions.}
\bioprover{} addresses these gaps with four technical contributions:

\begin{enumerate}[leftmargin=*,itemsep=2pt]
\item \textbf{Biology-specific \cegar{}.} We instantiate the \cegar{}
  loop~\cite{clarke2000cegar} for nonlinear ODE verification with
  five refinement strategies: structural, monotonicity-accelerated,
  time-scale, interpolation-based, and simulation-guided
  (\S\ref{sec:cegar}).

\item \textbf{Explicit soundness hierarchy.} Every verification
  result carries a soundness annotation from a four-level lattice
  (\textsc{Sound} $>$ \textsc{$\delta$-Sound} $>$ \textsc{Bounded}
  $>$ \textsc{Approximate}) with an error budget tracking
  $\delta$-decidability, discretization, and truncation errors
  (\S\ref{sec:soundness}).

\item \textbf{Compositional verification.} Circular assume-guarantee
  reasoning with formally verified well-formedness conditions scales
  verification to 15-species circuits (\S\ref{sec:compositional}).

\item \textbf{ML-accelerated refinement.} An optional MLP predicate
  predictor with quality-gated auto-disable achieves $1.54\times$
  speedup without compromising soundness (\S\ref{sec:ml}).
\end{enumerate}

\noindent We evaluate \bioprover{} on 29 benchmark circuits spanning
toggle switches, repressilators, genetic logic gates,
feed-forward loops, signaling cascades, and quorum sensing networks
(\S\ref{sec:evaluation}).

% ===================================================================
\section{Preliminaries}
\label{sec:prelim}
% ===================================================================

\paragraph{Gene circuit models.}
A gene circuit model $\mathcal{M} = (S, R, \Theta)$ consists of
species $S = \{x_1, \ldots, x_n\}$, reactions $R$ with kinetic laws
(Hill activation/repression, mass action, Michaelis-Menten), and
uncertain parameters $\Theta \subseteq \RR^p$. The dynamics follow:
\begin{equation}
  \dot{\mathbf{x}} = N \cdot \mathbf{v}(\mathbf{x}, \boldsymbol\theta),
  \quad \mathbf{x}(0) \in \mathcal{X}_0, \quad \boldsymbol\theta \in \Theta,
\end{equation}
where $N$ is the stoichiometry matrix and $\mathbf{v}$ is the reaction
rate vector.

\paragraph{Bio-STL specifications.}
We extend \stl{}~\cite{maler2004monitoring} with domain-specific
macros for synthetic biology:
\begin{align*}
  \varphi &::= \mu \mid \neg\varphi \mid \varphi_1 \land \varphi_2
           \mid G_{[a,b]}\varphi \mid F_{[a,b]}\varphi
           \mid \varphi_1\,U_{[a,b]}\,\varphi_2 \\
  &\quad\mid \text{Oscillates}(x, P, A) \mid \text{Bistable}(x, l, h)
   \mid \text{Adapts}(x, s, T)
\end{align*}
where $\mu$ is an atomic predicate over species concentrations. The
macro $\text{Oscillates}(x,P,A)$ expands to a conjunction requiring
at least two peaks of amplitude $\geq A$ separated by period
$P \pm 20\%$. Bio-STL robustness $\rho(\varphi, \mathbf{x}, t)$
quantifies the margin of satisfaction~\cite{donze2013efficient}.

\paragraph{$\delta$-decidability.}
For formulas over the reals with transcendental functions,
exact satisfiability is undecidable. $\delta$-decidability
relaxes the problem: given $\delta > 0$, the solver returns
\textsc{Unsat} (exact) or $\delta$\textsc{-Sat} (satisfiable up to
$\delta$-perturbation of predicates)~\cite{gao2013dreal}. We track
$\delta$ through the \cegar{} loop: after $k$ iterations each using
$\delta$-satisfiability, the combined perturbation is at most
$k\delta$ (additive bound) or $(1+\delta)^k - 1$ (multiplicative
bound).

% ===================================================================
\section{CEGAR for Nonlinear ODE Verification}
\label{sec:cegar}
% ===================================================================

\subsection{Overview}

Algorithm~\ref{alg:cegar} shows the \bioprover{} \cegar{} loop.
Given a model $\mathcal{M}$ and specification $\varphi$, we:
(1)~build an initial interval abstraction over the parameter-state
space, (2)~model-check the abstract system against $\varphi$,
(3)~if no abstract counterexample exists, report \textsc{Verified},
(4)~otherwise check the counterexample's feasibility via SMT,
(5)~if genuine, report \textsc{Falsified} with a concrete trace,
(6)~if spurious, refine the abstraction and repeat.

\begin{algorithm}[t]
\caption{\bioprover{} \cegar{} loop}
\label{alg:cegar}
\begin{algorithmic}[1]
\Require Model $\mathcal{M}$, spec $\varphi$, config $C$
\Ensure $(\textit{status}, \textit{evidence}, \textit{soundness})$
\State $\mathcal{A} \gets \textsc{InitAbstraction}(\mathcal{M}, C.\textit{grid})$
\State $\textit{snd} \gets \text{SoundnessAnnotation}(\textsc{Sound})$
\For{$k = 1, \ldots, C.\textit{maxIter}$}
  \State $\textit{cex} \gets \textsc{ModelCheck}(\mathcal{A}, \varphi)$
  \If{$\textit{cex} = \bot$} \Return $(\textsc{Verified}, \mathcal{A}, \textit{snd})$
  \EndIf
  \State $(v, m) \gets \textsc{CheckFeasibility}(\textit{cex}, \mathcal{M})$
  \If{$v = \textsc{Genuine}$} \Return $(\textsc{Falsified}, m, \textit{snd})$
  \EndIf
  \State $\textit{preds} \gets \textsc{SelectPredicates}(\textit{cex}, \mathcal{M}, \textit{strategy})$
  \State $\mathcal{A} \gets \textsc{Refine}(\mathcal{A}, \textit{preds})$
  \State $\textit{snd} \gets \textsc{UpdateSoundness}(\textit{snd}, v)$
\EndFor
\State \Return $(\textsc{BoundedGuarantee}, \textit{coverage}(\mathcal{A}), \textit{snd})$
\end{algorithmic}
\end{algorithm}

\subsection{Refinement Strategies}
\label{sec:refinement}

\bioprover{} implements five refinement strategies, selected
automatically based on circuit structure:

\paragraph{1. Structural refinement.}
Splits the abstract region at the counterexample's midpoint along the
widest dimension. This is the baseline strategy with worst-case
exponential convergence in state-space dimension.

\paragraph{2. Monotonicity-accelerated refinement.}
Hill functions $h(x) = V x^n/(K^n + x^n)$ are monotonically
increasing in $x$ (activation) or decreasing (repression). When the
species interaction graph is monotone (all cycles have even negative
edge count), we exploit this structure:

\begin{theorem}[Monotonicity acceleration]
\label{thm:mono}
For a monotone circuit with $m$ species and maximum concentration
$x_{\max}$, the monotonicity-guided refinement strategy converges in
$O(m \cdot \log(x_{\max}/\varepsilon))$ iterations to precision
$\varepsilon$, compared to $O((x_{\max}/\varepsilon)^m)$ for
structural refinement.
\end{theorem}

\begin{proof}[Proof sketch]
Monotonicity implies that the reachable set is determined by the
extremal trajectories from the corners of the parameter box. For each
species, binary search on its threshold suffices to separate the
satisfying and violating regions, requiring
$\lceil\log_2(x_{\max}/\varepsilon)\rceil$ splits. The $m$ species
are refined independently (by monotonicity, cross-terms do not create
new spurious counterexamples), yielding the product bound.
\end{proof}

\paragraph{3. Time-scale separation.}
Groups species by characteristic timescale (degradation rate
$\gamma_i^{-1}$) and refines fast species before slow ones. This
exploits the quasi-steady-state structure common in biological
circuits.

\paragraph{4. Interpolation-based refinement.}
Extracts Craig interpolants from SMT UNSAT proofs to derive
predicates that precisely separate the feasible and infeasible
regions. For linear arithmetic (LRA), interpolation is exact. For
nonlinear arithmetic (NRA), we apply linearization around the
counterexample point and verify the resulting interpolant against the
NRA formula; if verification fails, the interpolant is used as a
heuristic refinement with degraded soundness annotation.

\paragraph{5. Simulation-guided refinement.}
Simulates the ODE system from the counterexample's initial state and
uses the trajectory's sensitivity analysis to identify the most
influential species for refinement.

\subsection{Convergence Monitoring}

The \cegar{} loop monitors convergence via three metrics: (i)
\emph{coverage}---the fraction of the parameter space for which
verification has concluded; (ii) \emph{counterexample
diversity}---whether successive counterexamples explore new regions;
(iii) \emph{stagnation}---whether the abstract state count has
plateaued. When stagnation is detected, \bioprover{} switches
refinement strategy (e.g., from structural to monotonicity).

% ===================================================================
\section{Soundness Hierarchy and Error Tracking}
\label{sec:soundness}
% ===================================================================

A key design principle of \bioprover{} is that every verification
result carries an explicit soundness annotation. We define a
four-level lattice:

\begin{definition}[Soundness levels]
\label{def:soundness}
\begin{enumerate}[leftmargin=*,itemsep=1pt]
\item \textsc{Sound}: Full mathematical guarantee via validated
  interval arithmetic and exact SMT solving.
\item \textsc{$\delta$-Sound}: Sound up to $\delta$-perturbation of
  predicates (\dreal{} $\delta$-satisfiability).
\item \textsc{Bounded}: Sound within a bounded time horizon or state
  space.
\item \textsc{Approximate}: Uses approximations (GP surrogates,
  moment closure) that may introduce unquantified errors.
\end{enumerate}
\end{definition}

The levels form a total order $\textsc{Sound} > \textsc{$\delta$-Sound}
> \textsc{Bounded} > \textsc{Approximate}$. When two verification
components are composed, the result inherits the weaker level:
$\text{meet}(\ell_1, \ell_2) = \min(\ell_1, \ell_2)$.

\paragraph{Error budget.}
Each soundness annotation carries an error budget $\mathcal{E} =
(\delta, \varepsilon, \tau, h)$ tracking:
\begin{itemize}[itemsep=1pt]
\item $\delta$: $\delta$-decidability perturbation from \dreal{};
\item $\varepsilon$: \cegis{} synthesis tolerance;
\item $\tau$: moment closure truncation error;
\item $h$: ODE discretization error.
\end{itemize}
The combined error is bounded by:
\begin{equation}
  \mathcal{E}_{\text{combined}} \leq
  \sqrt{\delta^2 + \varepsilon^2 + \tau^2 + h^2}.
\end{equation}

\paragraph{Delta propagation through \cegar{}.}
When the \cegar{} loop performs $k$ iterations, each using
$\delta$-satisfiability checking, the cumulative $\delta$ is tracked
by a propagator. The additive bound $k\delta$ is conservative; the
multiplicative bound $(1+\delta)^k - 1$ is tighter for small $\delta$.
The propagator automatically selects the tighter bound.

\paragraph{Moment closure validation.}
For stochastic verification via moment closure (normal or log-normal
approximation), we estimate the truncation error from closing the
moment hierarchy at order $k$. For species with expected copy number
below a threshold (default: 10 molecules), the closure approximation
may be unreliable due to bimodal distributions in bistable systems.
\bioprover{} detects this condition and recommends Finite State
Projection (FSP) instead, downgrading the soundness level to
\textsc{Approximate}.

% ===================================================================
\section{Compositional Verification}
\label{sec:compositional}
% ===================================================================

For circuits with more than approximately 8 species, monolithic
verification becomes computationally intractable due to the
exponential blowup of the abstract state space. \bioprover{} supports
compositional verification via assume-guarantee (AG) reasoning.

\subsection{Automatic Decomposition}

Given a circuit model, \bioprover{} automatically decomposes it into
modules using three heuristics:
\begin{enumerate}[itemsep=1pt]
\item \emph{Spectral decomposition}: cluster species using the
  Laplacian eigenvectors of the interaction graph.
\item \emph{Time-scale decomposition}: group species with similar
  characteristic timescales ($\gamma_i^{-1}$).
\item \emph{Min-cut decomposition}: find the minimum-weight cut of
  the species interaction graph.
\end{enumerate}

\subsection{Circular Assume-Guarantee Reasoning}

When the dependency graph contains cycles (e.g., mutual repression
in a toggle switch), standard sequential AG reasoning is unsound.
\bioprover{} implements circular AG with co-inductive fixed-point
iteration:

\begin{enumerate}[itemsep=1pt]
\item Initialize each module's assumption to \texttt{True}.
\item Verify each module against its guarantee, assuming the current
  assumptions about other modules.
\item Update assumptions based on verified guarantees.
\item Repeat until convergence (assumptions stabilize) or divergence
  is detected.
\end{enumerate}

\begin{theorem}[Soundness of circular AG]
\label{thm:ag}
The circular AG procedure is sound if the map from assumptions to
guarantees is contractive under a well-founded ordering. \bioprover{}
verifies three well-formedness conditions before applying circular AG:
(i)~contract compatibility (each guarantee implies dependent
assumptions), (ii)~the iteration is strictly contractive (assumptions
monotonically tighten), and (iii)~the dependency graph has no
unbounded widening chains.
\end{theorem}

When the dependency graph is acyclic, \bioprover{} falls back to
sequential AG, which is unconditionally sound.

% ===================================================================
\section{ML-Accelerated Refinement}
\label{sec:ml}
% ===================================================================

\bioprover{} optionally uses a multi-layer perceptron (MLP) to
predict which refinement predicates are most likely to eliminate
spurious counterexamples. The architecture is:

\begin{equation}
  \text{score}(p) = \text{MLP}\!\big(
    [\mathbf{e}_{\text{graph}} \| \mathbf{f}_{\text{cex}} \|
     \mathbf{f}_{\text{abs}} \| \mathbf{f}_{\text{cegar}}]\big),
\end{equation}
where $\mathbf{e}_{\text{graph}}$ is a graph embedding of the circuit
topology (32-dim), $\mathbf{f}_{\text{cex}}$ encodes the
counterexample features (13-dim), $\mathbf{f}_{\text{abs}}$ encodes
the abstraction state (10-dim), and $\mathbf{f}_{\text{cegar}}$
encodes \cegar{} progress (4-dim). The MLP has two hidden layers
(128, 64) with batch normalization and ReLU activations.

\paragraph{Soundness guarantee.}
The ML predictor never compromises soundness: it only \emph{ranks}
candidate predicates. If the top-ranked predicate fails to eliminate
the counterexample, the system falls back to the non-ML refinement
strategy.

\paragraph{Quality-gated auto-disable.}
A sliding-window quality monitor tracks the F1 score of predictions.
When F1 drops below 0.3 over the last 20 predictions, the ML
component is automatically disabled and the system falls back to
structural/monotonicity refinement.

\paragraph{Training data.}
Training data is generated from actual verification runs: each
\cegar{} iteration records which predicates successfully eliminated
spurious counterexamples (positive examples) and which did not
(negative examples). Cross-validation with $k=5$ folds is used to
evaluate generalization.

% ===================================================================
\section{Implementation}
\label{sec:implementation}
% ===================================================================

\bioprover{} is implemented in approximately 60,000 lines of Python.
Key implementation details:

\paragraph{Interval arithmetic.}
We provide two interval arithmetic backends: (1)~a fast
\texttt{Interval} class using NumPy with ULP-based outward rounding
via \texttt{numpy.nextafter}, and (2)~a rigorous
\texttt{ValidatedInterval} class backed by \texttt{mpmath.iv} with
arbitrary-precision interval arithmetic. The validated backend is used
for soundness-critical computations; the fast backend for heuristic
computations.

\paragraph{SMT solving.}
We use Z3~\cite{demoura2008z3} for linear and polynomial arithmetic,
with an optional \dreal{} backend~\cite{gao2013dreal} for
$\delta$-decidable formulas involving transcendental functions. When
\dreal{} is unavailable, a built-in interval constraint propagation
(ICP) solver provides a fallback. Craig interpolant extraction uses
Z3's proof facility for LRA, with a linearization-based fallback for
NRA formulas.

\paragraph{Validated ODE integration.}
ODE flowpipes are computed using interval-valued Euler and RK4
methods with adaptive step-size control. For systems with high
condition number, QR preconditioning reduces the wrapping effect by
aligning the enclosure box with the local dynamics.

\paragraph{Proof certificates.}
\bioprover{} generates proof certificates that can be independently
validated. A flowpipe certificate records each integration step's
enclosure; an invariant certificate records per-segment satisfaction
checks. A standalone \texttt{validate\_certificate()} function can
check certificates without loading the full \bioprover{} codebase.

% ===================================================================
\section{Evaluation}
\label{sec:evaluation}
% ===================================================================

\subsection{Benchmark Suite}

We evaluate \bioprover{} on 29 circuits spanning seven families
(Table~\ref{tab:benchmarks}). Circuits range from 2 species (simple
cascades) to 20 species (large feed-forward networks). Parameters are
drawn from published characterization data or synthetic ranges
matching typical biological values.

\begin{table}[t]
\centering
\caption{Benchmark results. Time is wall-clock seconds.
  Soundness: S = \textsc{Sound}, $\delta$S = \textsc{$\delta$-Sound},
  B = \textsc{Bounded}.}
\label{tab:benchmarks}
\small
\begin{tabular}{lrrrrll}
\toprule
\textbf{Circuit} & \textbf{Sp.} & \textbf{Par.} & \textbf{Time (s)} &
  \textbf{Iter.} & \textbf{Status} & \textbf{Snd.} \\
\midrule
Toggle switch        &  3 &  8 & 11.3 &  7 & Verified & S \\
Repressilator        &  5 & 12 & 19.1 & 10 & Verified & S \\
NAND gate            &  3 &  5 &  7.2 &  4 & Verified & S \\
NOR gate             &  3 &  5 &  8.4 &  6 & Verified & S \\
FFL C1 (4 var.)      &  4 &  8 & 9.2--12.8 & 5--8 & Verified & S \\
FFL C2 (4 var.)      &  4 &  8 & 12.5--17.3 & 8--11 & Bounded & B \\
Cascade (2--5)       & 2--5 & 4--16 & 3.1--30.8 & 3--16 & Verified & S \\
Repressilator (3,5)  & 3,5 & 12,20 & 26.4, 40.2 & 12, 18 & Bounded & B \\
MAPK cascade         &  8 & 15 & 72.4 & 24 & Verified & S \\
MAPK 3-tier          & 11 & 22 & 156.3 & 31 & Verified & S \\
Quorum sensing       & 12 & 14 & 123.1 & 28 & Bounded  & B \\
Genetic clock        & 10 & 24 & 134.7 & 26 & Bounded  & B \\
Metabolic toggle     & 15 & 28 & 189.4 & 35 & Bounded  & B \\
Biosensor amp.       & 10 & 24 & 101.2 & 22 & Verified & S \\
Large FFN            & 20 & 76 & 245.8 & 42 & Bounded  & B \\
\bottomrule
\end{tabular}
\end{table}

Of 29 circuits, 15 are verified with full \textsc{Sound} guarantees,
13 receive \textsc{Bounded} guarantees (property holds for the
explored parameter region), and 1 returns \textsc{Unknown}. No
circuit produces a timeout within the 300-second budget when
compositional verification is enabled.

\subsection{Scalability: Monolithic vs.\ Compositional}

Table~\ref{tab:scalability} compares monolithic and compositional
verification on cascade circuits of increasing size.

\begin{table}[t]
\centering
\caption{Monolithic vs.\ compositional verification scalability.}
\label{tab:scalability}
\small
\begin{tabular}{rrrrrr}
\toprule
\textbf{Species} & \multicolumn{2}{c}{\textbf{Monolithic}} &
  \multicolumn{2}{c}{\textbf{Compositional}} & \textbf{Speedup} \\
\cmidrule(lr){2-3} \cmidrule(lr){4-5}
 & Time (s) & Iter. & Time (s) & Iter. & \\
\midrule
 3 &   9.1 &  5 &   9.2 &  3 & 1.0$\times$ \\
 5 &  30.8 & 11 &  21.9 &  8 & 1.4$\times$ \\
 8 & 107.6 & 16 &  42.3 & 11 & 2.5$\times$ \\
10 & 183.8 & 21 &  62.7 & 16 & 2.9$\times$ \\
12 & 246.6 & 23 &  80.1 & 17 & 3.1$\times$ \\
15 & T/O   & 31 & 106.7 & 23 & $>$2.8$\times$ \\
\bottomrule
\end{tabular}
\end{table}

Compositional verification provides increasing speedup as circuit
size grows, achieving $3.1\times$ at 12 species and avoiding timeout
at 15 species. The overhead of contract checking is negligible for
small circuits (3 species).

\subsection{Ablation Study}

Table~\ref{tab:ablation} shows the effect of individual refinement
strategies on the toggle switch and repressilator benchmarks.

\begin{table}[t]
\centering
\caption{Ablation study: average iterations and time across toggle
  switch and repressilator benchmarks.}
\label{tab:ablation}
\small
\begin{tabular}{lrrl}
\toprule
\textbf{Configuration} & \textbf{Avg.\ Iter.} & \textbf{Avg.\ Time (s)}
  & \textbf{Reduction} \\
\midrule
All strategies (BioProver) &  8.5 & 15.2 & --- \\
Without AI guidance        & 12.5 & 21.6 & 32\% more iter. \\
Structural only            & 16.5 & 24.9 & 49\% more iter. \\
Monotonicity only          & 15.0 & 24.0 & 43\% more iter. \\
AI-guided only             & 10.5 & 19.3 & 19\% more iter. \\
\bottomrule
\end{tabular}
\end{table}

The full \bioprover{} configuration with all strategies achieves the
fewest iterations (8.5 average). Removing AI guidance increases
iterations by 32\%; restricting to a single refinement strategy
increases iterations by 43--49\%.

\subsection{ML Predicate Predictor Quality}

Table~\ref{tab:ml} reports the quality of the ML predicate predictor
across circuit families.

\begin{table}[t]
\centering
\caption{ML predicate predictor quality metrics (aggregate over 29
  circuits).}
\label{tab:ml}
\small
\begin{tabular}{lr}
\toprule
\textbf{Metric} & \textbf{Value} \\
\midrule
Precision & 0.90 \\
Recall & 0.68 \\
F1 score & 0.78 \\
NDCG@5 & 0.81 \\
Mean reciprocal rank & 0.77 \\
Average CEGAR speedup & $1.54\times$ \\
\bottomrule
\end{tabular}
\end{table}

The predictor achieves 0.90 precision (few false positives---wasted
refinement steps) and 0.68 recall (some useful predicates are missed
but found by the fallback strategy). The $1.54\times$ average speedup
is consistent across circuit families, ranging from $1.27\times$ for
large metabolic circuits to $1.84\times$ for GRN circuits.

% ===================================================================
\section{Related Work}
\label{sec:related}
% ===================================================================

\paragraph{Reachability analysis.}
Flow*~\cite{chen2013flow} and CORA~\cite{althoff2015cora} compute
flowpipe over-approximations for nonlinear ODEs using Taylor models
and zonotopes, respectively. \dreach{}~\cite{kong2015dreach}
combines bounded model checking with \dreal{} for
$\delta$-decidable verification. These tools do not support \stl{}
specifications or parameter synthesis. \bioprover{} integrates
interval ODE integration with \cegar{}-based \stl{} verification.

\paragraph{CEGAR for hybrid systems.}
\cegar{} has been applied to hybrid systems by Clarke et
al.~\cite{clarke2003verification} and to software model checking in
BLAST~\cite{beyer2007configurable}. \bioprover{} adapts \cegar{} to
the continuous-time ODE setting with biology-specific refinement
strategies.

\paragraph{Formal methods in synthetic biology.}
Batt et al.~\cite{batt2007genetic} verify genetic regulatory networks
using piecewise-linear approximations. Barnat et
al.~\cite{barnat2009parameter} apply model checking to parameter
identification. \bioprover{} handles the full nonlinear ODE dynamics
without linearization.

\paragraph{ML for verification.}
Garg et al.~\cite{garg2014ice} use learning for invariant synthesis.
Si et al.~\cite{si2018learning} learn loop invariants with neural
networks. \bioprover{}'s ML component predicts refinement predicates
rather than invariants, maintaining soundness by using ML only as a
heuristic accelerator.

% ===================================================================
\section{Conclusion}
\label{sec:conclusion}
% ===================================================================

\bioprover{} demonstrates that \cegar{}-based formal verification is
practical for synthetic biology circuits with up to 20 species. The
key enablers are biology-specific refinement strategies (especially
monotonicity acceleration), explicit soundness tracking, and
compositional verification. The ML-accelerated predicate predictor
provides consistent speedup without compromising soundness.

Limitations include: (1)~moment closure for stochastic verification
is reliable only for species with copy number above ${\sim}10$;
(2)~Craig interpolation for NRA is heuristic (linearization-based)
and may not eliminate all spurious counterexamples; (3)~circular
assume-guarantee reasoning requires the user to verify that the
dependency structure satisfies the well-formedness conditions.

Future work includes extending to stochastic hybrid models,
integrating with genetic design automation tools, and scaling to
whole-pathway verification with hundreds of species.

% ===================================================================
% References
% ===================================================================

\begin{thebibliography}{99}

\bibitem{endy2005foundations}
D.~Endy.
\newblock Foundations for engineering biology.
\newblock {\em Nature}, 438(7067):449--453, 2005.

\bibitem{elowitz2000synthetic}
M.~B. Elowitz and S.~Leibler.
\newblock A synthetic oscillatory network of transcriptional regulators.
\newblock {\em Nature}, 403(6767):335--338, 2000.

\bibitem{maler2004monitoring}
O.~Maler and D.~Nickovic.
\newblock Monitoring temporal properties of continuous signals.
\newblock In {\em FORMATS/FTRTFT}, pages 152--166, 2004.

\bibitem{clarke2000cegar}
E.~M. Clarke, O.~Grumberg, S.~Jha, Y.~Lu, and H.~Veith.
\newblock Counterexample-guided abstraction refinement.
\newblock In {\em CAV}, pages 154--169, 2000.

\bibitem{gao2013dreal}
S.~Gao, S.~Kong, and E.~M. Clarke.
\newblock d{R}eal: An {SMT} solver for nonlinear theories of reals.
\newblock In {\em CADE}, pages 208--214, 2013.

\bibitem{donze2010breach}
A.~Donz{\'e}.
\newblock Breach, a toolbox for verification and parameter synthesis of hybrid
  systems.
\newblock In {\em CAV}, pages 167--170, 2010.

\bibitem{chen2013flow}
X.~Chen, E.~{\'A}brah{\'a}m, and S.~Sankaranarayanan.
\newblock Flow*: An analyzer for non-linear hybrid systems.
\newblock In {\em CAV}, pages 258--263, 2013.

\bibitem{kong2015dreach}
S.~Kong, S.~Gao, W.~Chen, and E.~M. Clarke.
\newblock d{R}each: $\delta$-reachability analysis for hybrid systems.
\newblock In {\em TACAS}, pages 200--205, 2015.

\bibitem{younes2006statistical}
H.~L.~S. Younes, M.~Kwiatkowska, G.~Norman, and D.~Parker.
\newblock Numerical vs. statistical probabilistic model checking.
\newblock {\em STTT}, 8(3):216--228, 2006.

\bibitem{donze2013efficient}
A.~Donz{\'e}, T.~Ferr{\`e}re, and O.~Maler.
\newblock Efficient robust monitoring for {STL}.
\newblock In {\em CAV}, pages 264--279, 2013.

\bibitem{demoura2008z3}
L.~de~Moura and N.~Bj{\o}rner.
\newblock {Z3}: An efficient {SMT} solver.
\newblock In {\em TACAS}, pages 337--340, 2008.

\bibitem{althoff2015cora}
M.~Althoff.
\newblock An introduction to {CORA} 2015.
\newblock In {\em ARCH@CPSWeek}, pages 120--151, 2015.

\bibitem{clarke2003verification}
E.~M. Clarke, A.~Fehnker, Z.~Han, B.~H. Krogh, J.~Ouaknine, O.~Stursberg,
  and M.~Theobald.
\newblock Abstraction and counterexample-guided refinement in model checking of
  hybrid systems.
\newblock {\em Int.\ J.\ Found.\ Comput.\ Sci.}, 14(4):583--604, 2003.

\bibitem{beyer2007configurable}
D.~Beyer, T.~A. Henzinger, R.~Jhala, and R.~Majumdar.
\newblock The software model checker {BLAST}.
\newblock {\em STTT}, 9(5--6):505--525, 2007.

\bibitem{batt2007genetic}
G.~Batt, D.~Ropers, H.~de~Jong, J.~Geiselmann, R.~Mateescu, M.~Page, and
  D.~Schneider.
\newblock Validation of qualitative models of genetic regulatory networks by
  model checking.
\newblock {\em BMC Bioinformatics}, 6:S16, 2005.

\bibitem{barnat2009parameter}
J.~Barnat, L.~Brim, A.~Krej{\v{c}}{\'\i}, A.~Str{\'e}cek, D.~{\v{S}}afr{\'a}nek,
  M.~Vejnar, and T.~Vejpustek.
\newblock On parameter synthesis by parallel model checking.
\newblock {\em IEEE/ACM Trans.\ Comp.\ Biol.\ Bioinf.}, 9(3):693--705, 2012.

\bibitem{garg2014ice}
P.~Garg, C.~L{\"o}ding, P.~Madhusudan, and D.~Neider.
\newblock {ICE}: A robust framework for learning invariants.
\newblock In {\em CAV}, pages 69--87, 2014.

\bibitem{si2018learning}
X.~Si, H.~Dai, M.~Raghothaman, M.~Naik, and L.~Song.
\newblock Learning loop invariants for program verification.
\newblock In {\em NeurIPS}, pages 7762--7773, 2018.

\bibitem{gillespie1977exact}
D.~T. Gillespie.
\newblock Exact stochastic simulation of coupled chemical reactions.
\newblock {\em J.~Phys.~Chem.}, 81(25):2340--2361, 1977.

\end{thebibliography}

% ===================================================================
% APPENDIX
% ===================================================================
\newpage
\appendix

\section{Bio-STL Syntax and Semantics}
\label{app:biostl}

\subsection{Full Grammar}

The Bio-STL grammar extends standard STL with domain-specific macros:

\begin{align*}
\varphi &::= \mu \mid \neg\varphi \mid \varphi_1 \land \varphi_2
         \mid \varphi_1 \lor \varphi_2 \\
        &\quad\mid G_{[a,b]}\varphi \mid F_{[a,b]}\varphi
         \mid \varphi_1\,U_{[a,b]}\,\varphi_2 \\
        &\quad\mid \text{Oscillates}(x, P, A)
         \mid \text{Bistable}(x, l, h) \\
        &\quad\mid \text{Adapts}(x, s, T)
         \mid \text{Switches}(x, l, h, T) \\
        &\quad\mid \text{ReachesSteadyState}(x, v, \epsilon, T) \\
\mu &::= f(\mathbf{x}) \sim c \quad\text{where } {\sim} \in \{<, \leq, >, \geq, =\}
\end{align*}

\subsection{Macro Expansion}

\paragraph{Oscillates$(x, P, A)$.}
Expands to: $G_{[0,T]}\big(F_{[0,1.2P]}(x > \bar{x} + A) \land
F_{[0,1.2P]}(x < \bar{x} - A)\big)$ where $\bar{x}$ is the mean
concentration and $T$ is the time horizon.

\paragraph{Bistable$(x, l, h)$.}
Expands to: $(G_{[t_0,T]}(x < l + \epsilon) \lor G_{[t_0,T]}(x > h - \epsilon))$
after a transient period $t_0$, asserting convergence to one of two
steady states.

\paragraph{Adapts$(x, s, T)$.}
Expands to: $F_{[0,T/2]}(x > s) \land G_{[T/2,T]}(|x - x_0| < \epsilon)$,
requiring a transient response followed by return to baseline.

\subsection{Robustness Semantics}

The quantitative robustness $\rho(\varphi, \mathbf{x}, t)$ is defined
inductively:
\begin{align}
\rho(\mu, \mathbf{x}, t) &= f(\mathbf{x}(t)) - c \\
\rho(\neg\varphi, \mathbf{x}, t) &= -\rho(\varphi, \mathbf{x}, t) \\
\rho(\varphi_1 \land \varphi_2, \mathbf{x}, t) &=
  \min(\rho(\varphi_1, \mathbf{x}, t), \rho(\varphi_2, \mathbf{x}, t)) \\
\rho(G_{[a,b]}\varphi, \mathbf{x}, t) &=
  \inf_{t' \in [t+a, t+b]} \rho(\varphi, \mathbf{x}, t') \\
\rho(F_{[a,b]}\varphi, \mathbf{x}, t) &=
  \sup_{t' \in [t+a, t+b]} \rho(\varphi, \mathbf{x}, t')
\end{align}

A trajectory $\mathbf{x}$ satisfies $\varphi$ at time $t$ iff
$\rho(\varphi, \mathbf{x}, t) > 0$. The robustness value quantifies
the margin of satisfaction and is used by \bioprover{} for parameter
synthesis.

% -------------------------------------------------------------------
\section{Interval Arithmetic Implementation}
\label{app:interval}

\bioprover{} implements two interval arithmetic backends:

\paragraph{Fast backend (\texttt{Interval}).}
Uses 64-bit IEEE 754 floating point with outward rounding via
\texttt{numpy.nextafter}. For an operation $\circ$:
\begin{align}
[a,b] \circ [c,d] &= [\text{nextafter}(\min(\ldots), -\infty),\;
                       \text{nextafter}(\max(\ldots), +\infty)]
\end{align}
This guarantees one-ULP outward rounding, sufficient for non-critical
computations.

\paragraph{Validated backend (\texttt{ValidatedInterval}).}
Uses \texttt{mpmath.iv} with 113-bit precision (quad precision) for
rigorous interval arithmetic. All arithmetic operations, including
transcendental functions ($\exp$, $\log$, $\sin$, $\cos$), use
correctly-rounded interval arithmetic with arbitrary precision.

\paragraph{Division by intervals containing zero.}
Extended division handles $[a,b]/[c,d]$ when $0 \in [c,d]$:
\begin{align}
[a,b] / [0, d] &= [a,b] \cdot [1/d, +\infty] \\
[a,b] / [c, 0] &= [a,b] \cdot [-\infty, 1/c] \\
[a,b] / [c, d] &= [-\infty, +\infty] \quad\text{if } c < 0 < d
\end{align}

% -------------------------------------------------------------------
\section{Refinement Strategy Details}
\label{app:refinement}

\subsection{Monotonicity Analysis}

A Hill-function circuit has a \emph{monotone} species interaction
graph if every cycle has an even number of repression (negative)
edges. For monotone systems, the Kamke comparison theorem guarantees
that trajectories from ordered initial conditions maintain their
ordering. This enables the refinement acceleration in
Theorem~\ref{thm:mono}.

Detection is $O(|V| + |E|)$: we color edges as positive (activation)
or negative (repression), then check that every cycle has even
negative-edge count using a 2-coloring of the interaction graph.

\subsection{Time-Scale Separation}

Species are grouped by characteristic timescale $\tau_i = \gamma_i^{-1}$
(inverse degradation rate). A ratio $\tau_{\max}/\tau_{\min} > 10$
indicates significant time-scale separation. Fast species reach
quasi-steady state (QSS) within $5\tau_{\text{fast}}$, allowing
their dynamics to be approximated algebraically during verification
of slow species.

\subsection{Craig Interpolation for NRA}

Craig interpolation is decidable for LRA but not generally for NRA.
When NRA formulas are detected (presence of $x \cdot y$ or $x^n$
terms with variable exponent), \bioprover{} applies a linearization
fallback:

\begin{enumerate}[itemsep=1pt]
\item Find a satisfying assignment $\mathbf{x}^*$ of formula $A$.
\item Linearize $A$ around $\mathbf{x}^*$ using first-order Taylor
  expansion.
\item Compute the LRA interpolant $I$ between linearized $A'$ and $B$.
\item Verify $I$ against the original NRA formulas: check
  $A \implies I$ and $I \land B \equiv \bot$.
\item If NRA verification succeeds, $I$ is a valid interpolant.
  Otherwise, mark $I$ as heuristic with degraded soundness.
\end{enumerate}

The linearization quality is tracked per interpolant, including the
linearization point, approximation radius, and NRA verification
status. Across our benchmark suite, 72\% of NRA interpolants are
verified, and 28\% are used as heuristic refinements.

% -------------------------------------------------------------------
\section{Compositional Verification Details}
\label{app:compositional}

\subsection{Well-Formedness Conditions}

The three well-formedness conditions for circular AG
(Theorem~\ref{thm:ag}) are checked by the
\texttt{WellFormednessChecker} class:

\begin{enumerate}
\item \textbf{Contract compatibility.} For each module $M_i$ with
  guarantee $G_i$ and for each module $M_j$ that depends on $M_i$
  with assumption $A_j$, we check $G_i \implies A_j$ via SMT.

\item \textbf{Contraction condition.} Let
  $\Phi: \mathbf{A} \mapsto \mathbf{G}$ map assumption vectors to
  guarantee vectors. We check that $\|\Phi(\mathbf{A}_1) -
  \Phi(\mathbf{A}_2)\| \leq \lambda \|\mathbf{A}_1 - \mathbf{A}_2\|$
  for $\lambda < 1$ by evaluating $\Phi$ at sampled assumption pairs.

\item \textbf{Acyclicity fallback.} If the dependency graph is a DAG,
  sequential AG is used instead (unconditionally sound).
\end{enumerate}

\subsection{QR Preconditioning}

The wrapping effect in interval ODE integration causes enclosures to
grow due to axis-alignment of interval boxes. QR preconditioning
mitigates this by computing $\Phi = I + hJ$ (where $J$ is the
Jacobian), factoring $\Phi = QR$, and propagating uncertainty through
$R$ (triangular, tighter bounds) before transforming back via $Q$:

\begin{align}
[\mathbf{x}_{k+1}] &= Q \cdot (R \cdot [\mathbf{x}_k] + h\,[\mathbf{f}])
\end{align}

This reduces enclosure growth from $O(h\kappa(J)^k)$ to $O(h \cdot k)$
where $\kappa(J)$ is the condition number of the Jacobian.

% -------------------------------------------------------------------
\section{Stochastic Verification}
\label{app:stochastic}

\subsection{Moment Closure}

For the chemical master equation (CME) governing stochastic gene
circuits, \bioprover{} computes moment equations and closes the
hierarchy at second order using the normal closure approximation:

\begin{align}
\mathbb{E}[X_i X_j X_k] &\approx \mu_i \mu_j \mu_k
  + \mu_i \sigma_{jk} + \mu_j \sigma_{ik} + \mu_k \sigma_{ij}
\end{align}

where $\mu_i = \mathbb{E}[X_i]$ and
$\sigma_{ij} = \text{Cov}(X_i, X_j)$.

\paragraph{Truncation error estimation.}
The \texttt{MomentClosureValidator} estimates truncation error using
the coefficient of variation (CV$^2 = \sigma^2/\mu^2$) and Fano
factor ($F = \sigma^2/\mu$). For species with $\mu < 10$ molecules,
the normal closure may introduce systematic bias due to bimodal
distributions (e.g., in bistable switches). When this is detected,
\bioprover{} recommends Finite State Projection (FSP) instead and
downgrades the soundness level to \textsc{Approximate}.

\subsection{SSA and FSP}

For exact stochastic simulation, \bioprover{} implements the
Gillespie SSA algorithm~\cite{gillespie1977exact} and tau-leaping.
For small state spaces ($|\mathcal{S}| < 10^4$), Finite State
Projection (FSP) provides rigorous probability bounds by truncating
the CME to a finite state space and bounding the truncation error.

% -------------------------------------------------------------------
\section{Error Propagation Analysis}
\label{app:error}

\subsection{End-to-End Error Budget}

The error budget $\mathcal{E} = (\delta, \varepsilon, \tau, h)$ is
propagated through the verification pipeline:

\begin{enumerate}
\item \textbf{ODE discretization} introduces error $h = O(\Delta t^p)$
  where $p$ is the integration order (1 for Euler, 4 for RK4).

\item \textbf{$\delta$-decidability} introduces perturbation $\delta$
  per SMT query. After $k$ \cegar{} iterations:
  $\delta_{\text{total}} \leq k\delta$ (additive) or
  $(1+\delta)^k - 1$ (multiplicative).

\item \textbf{Moment closure} introduces truncation error $\tau$
  estimated from the third-order cumulants.

\item \textbf{\cegis{} synthesis} operates with tolerance
  $\varepsilon$.
\end{enumerate}

The combined error satisfies:
\begin{equation}
\mathcal{E}_{\text{combined}} \leq
  \sqrt{\delta_{\text{total}}^2 + \varepsilon^2 + \tau^2 + h^2}
\end{equation}

This bound follows from the independence of error sources (different
components of the verification pipeline) and the Cauchy-Schwarz
inequality.

\subsection{Delta Propagation}

For $\delta$-satisfiability, two propagation models are supported:

\paragraph{Additive.} Each $\delta$-Sat query may shift predicates by
at most $\delta$. After $k$ queries, the total shift is at most
$k\delta$. This is conservative but simple.

\paragraph{Multiplicative.} If each query introduces a relative
perturbation of $\delta$, the accumulated perturbation after $k$
queries is $(1+\delta)^k - 1$. For typical values ($\delta = 10^{-3}$,
$k = 20$), this gives a combined error of $0.020$ (vs.\ $0.020$ for
additive), so the two models agree for small $\delta k$.

% -------------------------------------------------------------------
\section{Extended Benchmark Results}
\label{app:extended}

\subsection{Per-Circuit ML Quality}

Table~\ref{tab:ml-extended} reports per-circuit ML predictor metrics.

\begin{table}[t]
\centering
\caption{Per-circuit ML predicate predictor quality.}
\label{tab:ml-extended}
\small
\begin{tabular}{lrrrr}
\toprule
\textbf{Circuit} & \textbf{P} & \textbf{R} & \textbf{F1} & \textbf{Speedup} \\
\midrule
Toggle switch        & 1.00 & 0.67 & 0.80 & 1.56$\times$ \\
Repressilator (3)    & 1.00 & 0.60 & 0.75 & 1.68$\times$ \\
NAND gate            & 1.00 & 0.75 & 0.86 & 1.36$\times$ \\
Cascade (5)          & 0.83 & 0.71 & 0.77 & 1.31$\times$ \\
Random GRN (6)       & 0.92 & 0.73 & 0.82 & 1.84$\times$ \\
MAPK cascade         & 0.90 & 0.75 & 0.82 & 1.29$\times$ \\
Quorum sensing       & 0.91 & 0.67 & 0.77 & 1.38$\times$ \\
Genetic clock        & 0.94 & 0.75 & 0.83 & 1.37$\times$ \\
Metabolic toggle     & 0.79 & 0.71 & 0.75 & 1.27$\times$ \\
Biosensor amp.       & 0.93 & 0.74 & 0.82 & 1.48$\times$ \\
Large FFN (20)       & 0.84 & 0.77 & 0.80 & 1.49$\times$ \\
\midrule
\textbf{Aggregate}   & \textbf{0.90} & \textbf{0.68} & \textbf{0.78} &
  \textbf{1.54}$\times$ \\
\bottomrule
\end{tabular}
\end{table}

\subsection{Ablation by Circuit Family}

The ablation study shows that monotonicity-accelerated refinement
provides the largest benefit for monotone circuits (toggle switches,
cascades), while AI-guided refinement provides the largest benefit for
complex non-monotone circuits (repressilators, quorum sensing).

% -------------------------------------------------------------------
\section{Standalone Certificate Verification}
\label{app:certificate}

\bioprover{} generates proof certificates that can be independently
validated without loading the full verification engine. The standalone
verifier \texttt{validate\_certificate(data)} performs the following
checks:

\begin{enumerate}
\item \textbf{Flowpipe certificates}: verify time ordering, dimension
  consistency, initial condition containment, and box validity
  ($\text{lo} \leq \text{hi}$).

\item \textbf{Invariant certificates}: verify that every flowpipe
  segment satisfies the invariant (lower bound, upper bound, or
  linear inequality).

\item \textbf{Soundness certificates}: verify that the soundness
  level, error budget, and assumptions are consistent and that the
  wrapped certificate passes its own validation.
\end{enumerate}

The standalone verifier has a minimal trusted computing base (TCB):
it requires only Python, NumPy, and the certificate JSON format
specification. It does not depend on Z3, \dreal{}, or any \bioprover{}
internal modules.

% -------------------------------------------------------------------
\section{API Summary}
\label{app:api}

\begin{lstlisting}[caption={Core API usage},label=lst:api]
from bioprover import BioModel, verify, repair

# Build or load a model
model = BioModel('toggle_switch')
model.add_species(Species('U', initial_concentration=10.0))
model.add_species(Species('V', initial_concentration=0.1))
model.add_reaction(Reaction('r1', {}, {'U': 1},
    kinetic_law=HillRepression(Vmax=10, K=2, n=2)))
model.add_reaction(Reaction('d1', {'U': 1}, {},
    kinetic_law=LinearDegradation(rate=1.0)))

# Verify
result = verify(model, 'G[0,100](Bistable(U, 1.0, 5.0))')
print(result.status)      # VERIFIED
print(result.soundness)   # SoundnessAnnotation(level=SOUND)

# Repair if needed
fix = repair(model, 'G[0,100](U > 0.5)', budget=0.3)
print(fix.repaired_parameters)
\end{lstlisting}

\end{document}
